% !TeX root = ../../Book3.tex
% 纲要标题：### 13.6 Sobolev 映射空间与非线性算子
% 纲要位置：几何学计划纲要.md，第 2653 行。
% 正文状态：未撰写；本文件目前是框架。

% 纲要内容开始
% 
% 函数空间、微分阶数、系数正则性及参数范围同时固定；不对任意 Sobolev 指数无条件宣称复合算子光滑。
% 
% * 紧 $m$ 维源上的 $W^{k,p}$ 映射空间与 $kp>m$ 的连续嵌入
% * 指数图册、丛值截面及局部坐标转换
% * Sobolev 乘法、复合及微分损失的估计
% * 非线性微分算子的可微性和目标空间选择
% 
%   曲面上的 Cauchy–Riemann 模块固定 $p>2$，使用 $W^{1,p}\to L^p$ 的映射；高阶族与规范群作用另行核验所需指数，不将微分损失隐藏在“光滑作用”中。
% 
% * 椭圆自举、先验界与弱解的正则性恢复
% * 直接方法、约束临界点与规范对称的区别
% 
% ---
% 
% 纲要内容结束

\section{Sobolev 映射空间与非线性算子}
\label{g3:sec:1306}

待撰写

\subsection{\texorpdfstring{紧 \(m\) 维源上的 \(W^{k,p}\) 映射空间与 \(kp>m\) 的连续嵌入}{紧 m 维源上的 W\textasciicircum{}\{k,p\} 映射空间与 kp>m 的连续嵌入}}
\label{g3:subsec:1306:01}
待撰写

\subsection{指数图册、丛值截面及局部坐标转换}
\label{g3:subsec:1306:02}
待撰写

\subsection{Sobolev 乘法、复合及微分损失的估计}
\label{g3:subsec:1306:03}
待撰写

\subsection{非线性微分算子的可微性和目标空间选择}
\label{g3:subsec:1306:04}
待撰写

\subsection{椭圆自举、先验界与弱解的正则性恢复}
\label{g3:subsec:1306:05}
待撰写

\subsection{直接方法、约束临界点与规范对称的区别}
\label{g3:subsec:1306:06}
待撰写
% 纲要条目排为编号小节标题，并进入目录与书签；说明段仍保留于纲要注释中。
